\begin{table}
  \centering
  \scriptsize
  \begin{tabular}{p{3.1cm} p{2.6cm} p{5.6cm}}
    \toprule
    \textbf{维度} & \textbf{核心数字} & \textbf{解释} \\
    \midrule
    接受率（筛选） & 44\% vs. 51\%，下降约 7 个百分点（$p=0.007$） & 与“数字抵押减少逆向选择”一致: 一部分更可能违约的家庭在事前退出。 \\
    还款率 & 相对无担保贷款的 62\%，提高约 11 个百分点 & 数字抵押显著提高履约率，也提高贷款盈利能力。 \\
    全额偿还率 & 高出约 19 个百分点 & 不仅平均回款增加，右尾的“完全还清”也显著改善。 \\
    机制分解 & 约 $\frac{2}{3}$ 来自道德风险，约 $\frac{1}{3}$ 来自逆向选择 & 惊喜无担保组帮助把“筛选效应”和“激励效应”拆开。 \\
    异质性 & 道德风险主要集中在高风险家庭；逆向选择主要集中在低风险家庭 & 与模型关于不同类型借款人的预测一致。 \\
    \bottomrule
  \end{tabular}
\end{table}
